\subsection{Analytic functions}

3.2.1. Let U be an open subset of E and f a mapping from U into F. We say that f is of class \(C^{\omega}\) , or K-analytic (or simply analytic) in U if, for every point a in U, there exists a convergent series \(f_{a} \in \mathcal{H}(E; F)\) such that \(f(a + x) = f_{a}(x)\) for all x in E sufficiently close to zero. If K = R (resp. C), one also says that f is real analytic (resp. complex analytic or holomorphic). The analytic maps from U into F form a vector subspace, denoted \(\mathcal{C}^{\omega}(U; F)\) , of the space of all mappings from U into F.

For \(a \in \mathbf{U}\) , the formal series \(f_{a}\) is unique: it is called the power series expansion of \(f\) at the point \(a\) . If \(f_{a} = \sum_{\alpha}(f_{a})_{\alpha}\) (with \((f_{a})_{\alpha} \in \mathrm{P}_{\alpha}(\mathrm{E}_{1}, \ldots, \mathrm{E}_{n}; \mathrm{F})\) ), we set:

\[
\Delta^ {\alpha} f (a) = \left(f _ {a}\right) _ {\alpha}.
\]

3.2.2. If \(f \in \mathcal{C}^{\omega}(\mathrm{U};\mathrm{F})\) , the map \(\Delta^{\alpha}f: a \mapsto \Delta^{\alpha}f(a)\) from U into

\[
\mathbf {P} _ {\alpha} (\mathrm{E} _ {1}, \dots , \mathrm{E} _ {n}; \mathrm{F})
\]

is analytic. The map \(\Delta^{\alpha}: f \mapsto \Delta^{\alpha}f\) is a K-linear map from \(\mathcal{C}^{\omega}(\mathrm{U};\mathrm{F})\) to \(\mathcal{C}^{\omega}(\mathrm{U};\mathrm{P}_{\alpha}(\mathrm{E}_{1},\ldots,\mathrm{E}_{n};\mathrm{F}))\) . For \(a \in \mathrm{U}\) , we therefore have, for \(\alpha, \beta \in \mathbf{N}^{n}\) , \(\Delta^{\beta}(\Delta^{\alpha}f)(a) \in \mathrm{P}_{\beta}(\mathrm{E}_{1},\ldots,\mathrm{E}_{n};\mathrm{P}_{\alpha}(\mathrm{E}_{1},\ldots,\mathrm{E}_{n};\mathrm{F}))\) . If \(x = (x_{i}) \in \mathrm{E}\) , we therefore have \((\Delta^{\beta}(\Delta^{\alpha}f)(a))(x) \in \mathrm{P}_{\alpha}(\mathrm{E}_{1},\ldots,\mathrm{E}_{n};\mathrm{F})\) and \(((\Delta^{\beta}(\Delta^{\alpha}f)(a))(x))(x) \in \mathrm{F}\) . This element of F is equal to \(((\alpha, \beta)) (\Delta^{\alpha + \beta}f(a))(x)\) . This is expressed by writing:

\[
\Delta^ {\beta} \circ \Delta^ {\alpha} = ((\alpha , \beta)) \Delta^ {\alpha + \beta}.
\]

3.2.3. Strict convergence indicators (and strict radii of convergence when \(n=1\) ) of power series expansions of \(f\) and of \(\Delta^{\alpha}f\) at the same point \(a\) of U, are identical.

3.2.4. Let \(f \in \mathcal{C}^{\omega}(\mathrm{U};\mathrm{F})\) . Then \(f\) is strictly differentiable and infinitely differentiable in \(\mathbf{U}\) (if \(\mathbf{K} = \mathbf{R}\) , \(f\) is of class \(\mathbf{C}^{\infty}\) in \(\mathbf{U} \) ). The iterated derivatives of \(f\) are analytic, and their values at a point \(a\) are symmetric multilinear maps. One can then introduce the notation \(\mathbf{D}^{a}f\) for iterated partial derivatives as in no. 2.4.2. We have:

\[
\alpha ! \Delta^ {\alpha} f (a) (h) = D ^ {\alpha} f (a). (h, \dots , h)
\]

for all \(a \in \mathbf{U}\) and \(h \in \mathbf{E}\) , which we write:

\[
\alpha ! \Delta^ {\alpha} f = D ^ {\alpha} f.
\]

3.2.5. Let U be an open subset of E and f, g two analytic mappings from U into F. Let \(a \in U\) . For f and g to have contact of order \(\geqslant k\) at a, it is necessary and sufficient that \(\Delta^{\alpha}f(a) = \Delta^{\alpha}g(a)\) for all \(\alpha\) with \(|\alpha| \leqslant k\) . If f and g have contact of infinite order at a, they are equal in a neighborhood of a. The set of points of U where f and g have contact of infinite order is open and closed.

In particular, if U is connected and if there exists a nonempty open subset of U on which f and g are equal, then f = g ("principle of analytic continuation").

3.2.6. Let U be an open subset of E and let f be an analytic map from U into F. If the derivative Df of f is zero, then f is locally constant.

3.2.7. Suppose F is quasi-complete and let G be a complete normed space. Let g be an analytic map from an open subset U of E into G, and let f be an analytic map from an open subset V of G, containing g(U), into F. The composite map \(f \circ g\) is analytic in U. Suppose further that \(0 \in U\) and that \(g(0) = 0\) . The expansion of \(f \circ g\) as a power series at 0 is then obtained by substituting, in the expansion of f at 0, the power series expansion of g at 0 (3.1.9).

3.2.8. Let \(F_{1},\ldots,F_{m}\) be separated polynormed spaces and u a continuous multilinear mapping from \(F_{1}\times\cdots\times F_{m}\) into F. Let U be an open subset of E and \(f_{i}\in\mathcal{C}^{\omega}(U;F_{i})\) . The function \(u(f_{1},\ldots,f_{m})\) is analytic, and its power series expansion at a point \(a\in U\) is the series \(u((f_{1})_{a},\ldots,(f_{m})_{a})\) (3.1.8).

3.2.9. We assume \(F\) quasi-complete. Let \(f \in \mathcal{H}(E_1, \ldots, E_n; F)\); the function \(x \mapsto f(x)\) (3.1.7) is analytic in the open set \(C(f)\), the strict convergence domain of \(f\) . If \(n = 1\) and if \(\|a\| < \rho(f)\) , the strict radius of convergence of the expansion of \(f\) as a power series in \(a\) is at least equal to \(\rho(f) - \|a\|\) . If \(\rho(f) = +\infty\) , one says that \(f\) is an entire function.

3.2.10. Keep the hypotheses of 3.2.9. If K = C, the results of 3.2.9 remain exactly valid if one replaces C(f) by \(\tilde{C}(f)\) and \(\rho(f)\) by \(\tilde{\rho}(f)\) (for \(n = 1\) ). If K = R, the function \(x \mapsto f(x)\) is analytic in \(\tilde{C}(f)\) .

3.2.11. Suppose that \(E_{i} = K\) for \(1 \leqslant i \leqslant n\) . Let U be an open subset of \(K^{n}\) and \(f \in \mathcal{C}^{\omega}(K^{n}; F)\) . If \(0 \in U\) and if \(f_{0} = \sum_{\alpha} X^{\alpha} c_{\alpha}\) is the power series expansion of f at 0, the power series expansion of \(\Delta^{\alpha} f\) at 0 is written (after identifying \(P_{\alpha}(K^{n}; F)\) with F):

\[
(\Delta^ {\alpha} f) _ {0} = \sum_ {\beta} ((\alpha , \beta)) X ^ {\beta} c _ {\alpha + \beta}.
\]

In particular, for \(1 \leqslant i \leqslant n\) and for \(x \in \mathbf{C}(f_0)\) :

\[
\partial_ {i} f (x) = \sum_ {\alpha} (\alpha_ {i} + 1) x ^ {\alpha} c _ {\alpha + \varepsilon_ {i}}.
\]

